% GENERATED FILE. Edit canonical lessons, then run npm run generate:books.
% canonical-source-hash: fnv1a64:093ee9738cb7d541
% canonical-lessons: TE-C208-atuku, TE-C208-pukkilincu, TE-C208-canipovu, TE-C208-puttu, TE-C208-edurkonu

\chapter{To stick, To gargle, To die, To be born, To confront}
\label{ch:te-to-stick-to-gargle-to-die-to-be-born-to-confront}
\hlchaptermodality{208}

\begin{chapteropening}
\textbf{By the end of this chapter:} \emph{I can say to stick, to gargle, to die, to be born and to confront.}

Complete the last lesson of chapter 208: I can say to stick, to gargle, to die, to be born and to confront.
\end{chapteropening}

\section[\texorpdfstring{\te{అతుకు} (atuku)}{అతుకు (atuku)}]{\te{అతుకు} (atuku) — to stick, to be joined together}

\label{lesson:TE-C208-atuku}



\begin{quote}
Before the new one: say the Telugu for to till, then the Telugu for to grow.
\end{quote}

\subsection*{You'll want to know: \te{అతుకు}}
\textbf{\te{అతుకు}} — \emph{atuku} — \textquotedblleft{}to stick, to be joined together\textquotedblright{}.

\begin{grammarlens}[title={What you've built}]
One more word for this chapter.
\end{grammarlens}

\subsection*{Your turn}
\begin{itemize}
\raggedright
  \item \emph{Say it:} \emph{atuku}
  \item \emph{Say it:} \emph{atuku}, once more
  \item \emph{Say it:} say \emph{atuku}
  \item \emph{Recall:} say the Telugu for to till, then the Telugu for to grow, then say \emph{atuku} again
\end{itemize}

\begin{tcolorbox}[breakable,colback=teal!4,colframe=teal!35!black,title={Before you move on}]
What does \te{అతుకు} mean? (\textquotedblleft{}To stick, to be joined together\textquotedblright{}.) Say it once more.
\end{tcolorbox}
\section[\texorpdfstring{\te{పుక్కిలించు} (pukkiliñcu)}{పుక్కిలించు (pukkiliñcu)}]{\te{పుక్కిలించు} (pukkiliñcu) — to gargle}

\label{lesson:TE-C208-pukkilincu}



\begin{quote}
Before the new one: say the Telugu for to grow, then the Telugu for to stick.
\end{quote}

\subsection*{You'll want to know: \te{పుక్కిలించు}}
\textbf{\te{పుక్కిలించు}} — \emph{pukkiliñcu} — \textquotedblleft{}to gargle\textquotedblright{}.

\begin{grammarlens}[title={What you've built}]
One more word for this chapter.
\end{grammarlens}

\subsection*{Your turn}
\begin{itemize}
\raggedright
  \item \emph{Say it:} \emph{pukkiliñcu}
  \item \emph{Say it:} \emph{pukkiliñcu}, once more
  \item \emph{Say it:} say \emph{pukkiliñcu}
  \item \emph{Recall:} say the Telugu for to grow, then the Telugu for to stick, then say \emph{pukkiliñcu} again
\end{itemize}

\begin{tcolorbox}[breakable,colback=teal!4,colframe=teal!35!black,title={Before you move on}]
What does \te{పుక్కిలించు} mean? (\textquotedblleft{}To gargle\textquotedblright{}.) Say it once more.
\end{tcolorbox}
\section[\texorpdfstring{\te{చనిపోవు} (canipōvu)}{చనిపోవు (canipōvu)}]{\te{చనిపోవు} (canipōvu) — to die}

\label{lesson:TE-C208-canipovu}



\begin{quote}
Before the new one: say the Telugu for to stick, then the Telugu for to gargle.
\end{quote}

\subsection*{You'll want to know: \te{చనిపోవు}}
\textbf{\te{చనిపోవు}} — \emph{canipōvu} — \textquotedblleft{}to die\textquotedblright{}.

\begin{grammarlens}[title={What you've built}]
One more word for this chapter.
\end{grammarlens}

\subsection*{Your turn}
\begin{itemize}
\raggedright
  \item \emph{Say it:} \emph{canipōvu}
  \item \emph{Say it:} \emph{canipōvu}, once more
  \item \emph{Say it:} say \emph{canipōvu}
  \item \emph{Recall:} say the Telugu for to stick, then the Telugu for to gargle, then say \emph{canipōvu} again
\end{itemize}

\begin{tcolorbox}[breakable,colback=teal!4,colframe=teal!35!black,title={Before you move on}]
What does \te{చనిపోవు} mean? (\textquotedblleft{}To die\textquotedblright{}.) Say it once more.
\end{tcolorbox}
\section[\texorpdfstring{\te{పుట్టు} (puṭṭu)}{పుట్టు (puṭṭu)}]{\te{పుట్టు} (puṭṭu) — to be born}

\label{lesson:TE-C208-puttu}



\begin{quote}
Before the new one: say the Telugu for to gargle, then the Telugu for to die.
\end{quote}

\subsection*{You'll want to know: \te{పుట్టు}}
\textbf{\te{పుట్టు}} — \emph{puṭṭu} — \textquotedblleft{}to be born\textquotedblright{}.

\begin{grammarlens}[title={What you've built}]
One more word for this chapter.
\end{grammarlens}

\subsection*{Your turn}
\begin{itemize}
\raggedright
  \item \emph{Say it:} \emph{puṭṭu}
  \item \emph{Say it:} \emph{puṭṭu}, once more
  \item \emph{Say it:} say \emph{puṭṭu}
  \item \emph{Recall:} say the Telugu for to gargle, then the Telugu for to die, then say \emph{puṭṭu} again
\end{itemize}

\begin{tcolorbox}[breakable,colback=teal!4,colframe=teal!35!black,title={Before you move on}]
What does \te{పుట్టు} mean? (\textquotedblleft{}To be born\textquotedblright{}.) Say it once more.
\end{tcolorbox}
\section[\texorpdfstring{\te{ఎదుర్కొను} (edurkonu)}{ఎదుర్కొను (edurkonu)}]{\te{ఎదుర్కొను} (edurkonu) — to confront, to face}

\label{lesson:TE-C208-edurkonu}



\begin{quote}
Before the new one: say the Telugu for to die, then the Telugu for to be born.
\end{quote}

\subsection*{You'll want to know: \te{ఎదుర్కొను}}
\textbf{\te{ఎదుర్కొను}} — \emph{edurkonu} — \textquotedblleft{}to confront, to face\textquotedblright{}.

\begin{grammarlens}[title={What you've built}]
One more word for this chapter.
\end{grammarlens}

\subsection*{Your turn}
\begin{itemize}
\raggedright
  \item \emph{Say it:} \emph{edurkonu}
  \item \emph{Say it:} \emph{edurkonu}, once more
  \item \emph{Say it:} say \emph{edurkonu}
  \item \emph{Recall:} say the Telugu for to die, then the Telugu for to be born, then say \emph{edurkonu} again
\end{itemize}

\begin{tcolorbox}[breakable,colback=teal!4,colframe=teal!35!black,title={Before you move on}]
What does \te{ఎదుర్కొను} mean? (\textquotedblleft{}To confront, to face\textquotedblright{}.) Say it once more.
\end{tcolorbox}
